\chapter{Sound and Complete Solving for Multi-Width Parametric Bitvectors}
\label{ch:multi-width-bv}

\paragraph*{Attribution.} This chapter is based on \emph{Sound and Complete
Solving for Multi-Width Parametric Bitvectors via Principled
Reductions}~\citep{bhat2026multiwidth}, conditionally accepted at OOPSLA 2026.

\lettrine{F}{ixed-width} bitvector solvers (\qfbv{}) are fast and ubiquitous,
and \cref{ch:single-width-bv} showed how to decide predicates over a single
\emph{symbolic} width. However, the theory of \emph{parametric} bitvectors
(PBV), where widths are symbolic, is much less well understood in its full
generality. The theory of \emph{multi-width} PBV, where expressions may
involve $n$ distinct symbolic widths ($\text{PBV}_n$), is particularly
challenging: the only previously existing complete approach for bounded PBV
(where all widths have a concrete upper bound) is exhaustive enumeration,
requiring one \qfbv{} solver call for each of the exponentially many width
assignments. This is a significant bottleneck in tools, such as
Alive2~\citep{lopes2021alive2} and \hydra{}~\citep{mukherjee2024hydra}, that
formally reason about compiler optimizations.

To address this problem, we first prove that any $\text{PBV}_n$ formula can be
reduced to an equisatisfiable \emph{mono-width} ($\text{PBV}_1$) formula with
only a linear increase in formula size. The key idea is to encode symbolic
widths as \emph{bitmasks}. This reduction lets us create two solvers for
flavors of multi-width PBV:
\begin{enumerate}
  \item a sound and complete \emph{bounded} PBV solver, which instantiates
  the width variable of the $\text{PBV}_1$ formula to a concrete bound and
  therefore requires only a \emph{single} \qfbv{} call---in practice proving
  LLVM rewrites in seconds that enumeration fails to prove in hours; and
  \item by composing the reduction with the automata-theoretic decision
  procedures of \cref{ch:single-width-bv}, a new \emph{sound and complete
  decision procedure} for a fragment of $\text{PBV}_n$ with parametric
  widths, subsuming the prior state-of-the-art fragment of linear and bitwise
  operations by adding zero and sign extension.
\end{enumerate}
All our solvers are implemented in Lean, with mechanized proofs of soundness
and completeness for the unbounded solver. Empirically, the equisatisfiable
reduction from $\text{PBV}_n$ to $\text{PBV}_1$ turns exponential enumeration
into a single \qfbv{} query that nearly saturates standard PBV benchmarks,
while our unbounded solvers prove several times as many all-width problems as
the state-of-the-art CVC5-based solver.

\section{The Algorithmic Idea: Masks and Binary Values}
\label{sec:mwbv:idea}

\sid{Port from papers/paper-multi-width-bv/paper.tex \S 2 (lines 671--826).}

\section{A Unified Presentation of Zero, One, Many Width Parametric Bitvector Theory}
\label{sec:mwbv:theory}

\sid{Port from papers/paper-multi-width-bv/paper.tex \S 3 (lines 827--1034).}

\section{Equisatisfiable Reduction: Multi-Width to Mono-Width}
\label{sec:mwbv:reduction}

\sid{Port from papers/paper-multi-width-bv/paper.tex \S 4 (lines 1035--1260).}

\section{Extending Automata-Based PBV Solving to Multiple Widths}
\label{sec:mwbv:automata}

\sid{Port from papers/paper-multi-width-bv/paper.tex \S 5 (lines 1261--1764).}

\section{Solver Implementation and Mechanization}
\label{sec:mwbv:implementation}

\sid{Port from papers/paper-multi-width-bv/paper.tex \S 6 (lines 1765--1961).}

\section{Evaluation}
\label{sec:mwbv:evaluation}

\sid{Port from papers/paper-multi-width-bv/paper.tex \S 7 (lines 1962--2236).
Plots live under chapters/multi-width-bv/plots/output/<experiment>/; the
generated numbers (cactus.tex, stats.tex) were copied with commands prefixed
Mw (e.g. \textbackslash MwMachineHostname) --- \textbackslash input only the
experiments the text uses, e.g.
\textbackslash input\{chapters/multi-width-bv/plots/output/alive-cactus/cactus.tex\}.}

\section{Related Work}
\label{sec:mwbv:related}

\sid{Port from papers/paper-multi-width-bv/paper.tex \S 8 (lines 2237--2266).}

\section{Conclusions}
\label{sec:mwbv:conclusion}

\sid{Port from papers/paper-multi-width-bv/paper.tex \S 9 (lines 2267--),
and add a hand-off paragraph: from richer widths to richer values,
motivating \cref{ch:floating-point}.}
